% Part: first-order-logic
% Chapter: syntax-and-semantics
% Section: main-operator

\documentclass[../../../include/open-logic-section]{subfiles}

\begin{document}

\olfileid{fol}{syn}{mai}

\olsection{公式的\printtoken{S}{main operator}}

\begin{explain}
在讨论公式~$!A$ 的构造时，我们常常需要提及最后使用的那个算子。
这个算子称为~$!A$ 的\emph{主联结词}。
直观上看，它就是~$!A$ 中“最外层”的算子。
例如，$\lnot !A$ 的主联结词是~$\lnot$，
$(!A \lor !B)$ 的主联结词是~$\lor$，等等。
\end{explain}

\begin{defn}[主联结词]
\ollabel{def:main-op}
公式~$!A$ 的\emph{主联结词}定义如下：
\begin{enumerate}
\item \indcase*{!A}{!A}{$\indfrm$ 没有主联结词。}

\tagitem{prvNot}{\indcase{!A}{\lnot !B}{$\indfrm$ 的主联结词是~$\lnot$。}}{}

\tagitem{prvAnd}{\indcase{!A}{(!B \land !C)}{$\indfrm$ 的主联结词是~$\land$。}}{}

\tagitem{prvOr}{\indcase{!A}{(!B \lor !C)}{$\indfrm$ 的主联结词是~$\lor$。}}{}

\tagitem{prvIf}{\indcase{!A}{(!B \lif !C)}{$\indfrm$ 的主联结词是~$\lif$。}}{}

\tagitem{prvIff}{\indcase{!A}{(!B \liff !C)}{$\indfrm$ 的主联结词是~$\liff$。}}{}

\tagitem{prvAll}{\indcase{!A}{\lforall[x][!B]}{$\indfrm$ 的主联结词是~$\lforall$。}}{}

\tagitem{prvEx}{\indcase{!A}{\lexists[x][!B]}{$\indfrm$ 的主联结词是~$\lexists$。}}{}
\end{enumerate}
\end{defn}

在每种情形下，我们指的是主联结词在公式中被指明的那个特定\emph{出现}。
例如，由于公式
$((!D \lif !E) \lif (!E \lif !D))$
具有形式 $(!B \lif !C)$，其中 $!B$ 是 $(!D \lif !E)$ 而 $!C$ 是 $(!E \lif !D)$，那么第二次出现的 $\lif$ 就是主联结词。

\begin{explain}
这是一个函数的\emph{递归}定义，该函数将所有非原子公式映射到其主联结词的出现。
由于公式是归纳定义的，每个公式~$!A$ 都满足 \olref{def:main-op} 中的一种情形。
这保证了每个非原子公式~$!A$ 都有一个主联结词。
又因为每个公式只满足这些条件中的一个，并且在各种情形下构造~$!A$ 所用的较小公式都是唯一确定的，因此~$!A$ 的主联结词出现也是唯一的，所以我们确实定义了一个函数。
\end{explain}

我们根据主联结词的符号，按照 \olref{tab:main-op} 中的名称来称呼公式。
\iftag{defNot,defOr,defAnd,defIf,defIff,defTrue,defFalse,defEx,defAll}
{然而请记住，被定义的算子并不严格出现在公式中。它们只是缩写，因此严格地说，它们不能作为公式的主联结词。所以在关于所有公式的证明中，无需对它们单独处理。}

\begin{table}[!h]
\centering
\begin{tabular}{c | c | c}
主联结词 & 公式类型 & 示例\\
\hline
无 & 原子（公式） &
\iftag{prvFalse}{$\lfalse$,}{}
\iftag{prvTrue}{$\ltrue$,}{}
$\Atom{R}{t_1, \dots, t_n}$,
$\eq[t_1][t_2]$\\
$\lnot$ & 否定 & $\lnot !A$ \\
$\land$ & 合取 & $(!A \land !B$) \\
$\lor$ & 析取 & $(!A \lor !B$) \\
$\lif$ & 条件句 & $(!A \lif !B$) \\
$\liff$ & 双条件句 & $(!A \liff !B)$ \\
$\lforall[][]$ & 全称（公式）& $\lforall[x][!A]$ \\
$\lexists[][]$ & 存在（公式）& $\lexists[x][!A]$
\end{tabular}
\caption{公式的主联结词与名称}
\ollabel{tab:main-op}
\end{table}

\end{document}